% GENERATED FILE. Edit canonical lessons, then run npm run generate:books.
% canonical-source-hash: fnv1a64:a54f0ad3237459ce
% canonical-lessons: TA-C293-poruttatu, TA-C293-vantiruntal, TA-W293-six-lines, TA-R293-depends-and-would

\chapter{It Depends and What Would Have Been}
\label{ch:ta-it-depends-and-what-would-have-been}
\hlchaptermodality{293}

\begin{chapteropening}
\textbf{By the end of this chapter:} \emph{I can say what something depends on with … -aip poṟuttatu, say what would have happened with -iruntāl … -iruppēṉ, and write six connected lines about a trip.}

Complete the last lesson of chapter 293: say what something depends on and what would have been, after writing six connected lines about a trip.
\end{chapteropening}

\section[\texorpdfstring{-\ta{ஐப்} \ta{பொறுத்தது} (-aip poṟuttatu)}{-ஐப் பொறுத்தது (-aip poṟuttatu)}]{-\ta{ஐப்} \ta{பொறுத்தது} (-aip poṟuttatu) — it depends}

\label{lesson:TA-C293-poruttatu}



\begin{quote}
Say \textquotedblleft{}If it rains, I will stay at home.\textquotedblright{} (\emph{maḻai peytāl nāṉ vīṭṭil iruppēṉ}.) Then the word for the weather. (\emph{vāṉilai}.)
\end{quote}

\subsection*{You'll want to know: … -\ta{ஐப்} \ta{பொறுத்தது}}
\begin{quote}
\textbf{\ta{நாளை} \ta{வருவீர்களா}? \ta{அது} \ta{வானிலையைப்} \ta{பொறுத்தது}.}
\end{quote}

\begin{quote}
\emph{nāḷai varuvīrkaḷā? atu vāṉilaiyaip poṟuttatu.} — \textquotedblleft{}Will you come tomorrow? It depends on the weather.\textquotedblright{}
\end{quote}

\textbf{\ta{வருவீர்களா}} (\emph{varuvīrkaḷā}, \textquotedblleft{}will you come?\textquotedblright{}). \textbf{\ta{பொறுத்தது}} is \textquotedblleft{}it rests on\textquotedblright{}, from \textbf{\ta{பொறு}}, \textquotedblleft{}to bear\textquotedblright{}, the verb inside \textbf{\ta{பொறுத்துக்கொள்}}. The thing it depends on takes \textbf{-\ta{ஐ}}, and the first letter of \textbf{\ta{பொறுத்தது}} doubles in writing: \textbf{\ta{வானிலையைப்} \ta{பொறுத்தது}}. With a person, it hands them the choice: \textbf{\ta{அது} \ta{உங்களைப்} \ta{பொறுத்தது}} (\emph{atu uṅkaḷaip poṟuttatu}, \textquotedblleft{}it's up to you\textquotedblright{}).

Before another verb it becomes \textbf{\ta{பொறுத்து}}, \textquotedblleft{}depending on\textquotedblright{}:

\begin{quote}
\textbf{\ta{விலையைப்} \ta{பொறுத்து} \ta{வாங்குவேன்}.}
\end{quote}

\begin{quote}
\emph{vilaiyaip poṟuttu vāṅkuvēṉ.} — \textquotedblleft{}I'll buy it depending on the price.\textquotedblright{}
\end{quote}

A more formal partner, for writing: \textbf{… -\ta{ஐச்} \ta{சார்ந்திருக்கிறது}} (\emph{… -aic cārntirukkiṟatu}, \textquotedblleft{}it relies on …\textquotedblright{}), as in \textbf{\ta{இது} \ta{உங்கள்} \ta{முடிவைச்} \ta{சார்ந்திருக்கிறது}} (\emph{itu uṅkaḷ muṭivaic cārntirukkiṟatu}, \textquotedblleft{}this depends on your decision\textquotedblright{}).

\subsection*{Your turn}
Say these aloud:

\begin{itemize}
\raggedright
  \item \textquotedblleft{}Will you come tomorrow? It depends on the weather.\textquotedblright{} — \emph{nāḷai varuvīrkaḷā? atu vāṉilaiyaip poṟuttatu.}
  \item \textquotedblleft{}It's up to you\textquotedblright{} — \emph{atu uṅkaḷaip poṟuttatu}
  \item \textquotedblleft{}I'll buy it depending on the price\textquotedblright{} — \emph{vilaiyaip poṟuttu vāṅkuvēṉ}
  \item one plan of yours that depends on something, and what it depends on
\end{itemize}

\begin{tcolorbox}[breakable,colback=teal!4,colframe=teal!35!black,title={Before you move on}]
\textquotedblleft{}It depends on the weather\textquotedblright{}? (\textbf{\ta{வானிலையைப்} \ta{பொறுத்தது}}.) \textquotedblleft{}It's up to you\textquotedblright{}? (\textbf{\ta{அது} \ta{உங்களைப்} \ta{பொறுத்தது}}.) \textquotedblleft{}Depending on the price\textquotedblright{}? (\textbf{\ta{விலையைப்} \ta{பொறுத்து}}.)
\end{tcolorbox}
\section[\texorpdfstring{\ta{வந்திருந்தால்} (vantiruntāl)}{வந்திருந்தால் (vantiruntāl)}]{\ta{வந்திருந்தால்} (vantiruntāl) — if it had, it would have}

\label{lesson:TA-C293-vantiruntal}



\begin{quote}
Say \textquotedblleft{}having come\textquotedblright{} (\emph{vantu}) and \textquotedblleft{}if (it) is\textquotedblright{} (\emph{iruntāl}). Then \textquotedblleft{}It depends on the weather.\textquotedblright{} (\emph{atu vāṉilaiyaip poṟuttatu}.)
\end{quote}

\subsection*{You'll want to know: -\ta{இருந்தால்} … -\ta{இருப்பேன்}}
\begin{quote}
\textbf{\ta{எனக்கு} \ta{நேரம்} \ta{இருந்திருந்தால்}, \ta{நான்} \ta{வந்திருப்பேன்}.}
\end{quote}

\begin{quote}
\emph{eṉakku nēram iruntiruntāl, nāṉ vantiruppēṉ.} — \textquotedblleft{}If I had had time, I would have come.\textquotedblright{}
\end{quote}

Two pieces you know make the if-clause: the \textbf{-\ta{உ}} form and \textbf{\ta{இருந்தால்}}, \textquotedblleft{}if (it) is\textquotedblright{}. \textbf{\ta{வந்து}} + \textbf{\ta{இருந்தால்}} $\to$ \textbf{\ta{வந்திருந்தால்}}, \textquotedblleft{}if (someone) had come\textquotedblright{}; the \textbf{-\ta{உ}} drops before the \textbf{\ta{இ}}. With \textbf{\ta{இரு}} itself: \textbf{\ta{இருந்திருந்தால்}}, \textquotedblleft{}if (there) had been\textquotedblright{}.

The main clause takes \textbf{\ta{இருப்பேன்}}, \textquotedblleft{}I will be\textquotedblright{}, on the same \textbf{-\ta{உ}} form: \textbf{\ta{வந்திருப்பேன்}}, \textquotedblleft{}I would have come\textquotedblright{}. After \textbf{\ta{போய்}} the \textbf{\ta{இ}} stays whole: \textbf{\ta{போயிருப்பேன்}}, \textquotedblleft{}I would have gone\textquotedblright{}.

Both halves are about what did \textbf{not} happen: there was no time, and nobody came.

\subsection*{You'll want to know: would have, for everyone}
\begin{quote}
\textbf{\ta{நீங்கள்} \ta{வந்திருந்தால்}, \ta{நாம்} \ta{சேர்ந்து} \ta{கடலுக்குப்} \ta{போயிருப்போம்}.}
\end{quote}

\begin{quote}
\emph{nīṅkaḷ vantiruntāl, nām cērntu kaṭalukkup pōyiruppōm.} — \textquotedblleft{}If you had come, we would have gone to the sea together.\textquotedblright{}
\end{quote}

The ending of \textbf{\ta{இருப்பேன்}} changes for the person, as every future does: \textbf{\ta{போயிருப்போம்}}, \textquotedblleft{}we would have gone\textquotedblright{}; \textbf{\ta{வந்திருப்பான்}}, \textquotedblleft{}he would have come\textquotedblright{}. For \textquotedblleft{}it\textquotedblright{}, \textbf{\ta{இருக்கும்}}:

\begin{quote}
\textbf{\ta{அது} \ta{நன்றாக} \ta{இருந்திருக்கும்}.}
\end{quote}

\begin{quote}
\emph{atu naṉṟāka iruntirukkum.} — \textquotedblleft{}That would have been nice.\textquotedblright{}
\end{quote}

The same form also makes a guess about the past. \textbf{\ta{அவன்} \ta{வந்திருப்பான்}} can mean \textquotedblleft{}he must have come (by now)\textquotedblright{}. The rest of the sentence tells you which.

\subsection*{Your turn}
Say these aloud:

\begin{itemize}
\raggedright
  \item \textquotedblleft{}If I had had time, I would have come\textquotedblright{} — \emph{eṉakku nēram iruntiruntāl, nāṉ vantiruppēṉ}
  \item \textquotedblleft{}If you had come, we would have gone to the sea together\textquotedblright{} — \emph{nīṅkaḷ vantiruntāl, nām cērntu kaṭalukkup pōyiruppōm}
  \item \textquotedblleft{}That would have been nice\textquotedblright{} — \emph{atu naṉṟāka iruntirukkum}
\end{itemize}

\begin{tcolorbox}[breakable,colback=teal!4,colframe=teal!35!black,title={Before you move on}]
\textquotedblleft{}If (you) had come\textquotedblright{}? (\textbf{\ta{வந்திருந்தால்}}.) \textquotedblleft{}I would have come\textquotedblright{}? (\textbf{\ta{வந்திருப்பேன்}}.) \textquotedblleft{}I would have gone\textquotedblright{}? (\textbf{\ta{போயிருப்பேன்}}.) \textquotedblleft{}It would have been nice\textquotedblright{}? (\textbf{\ta{நன்றாக} \ta{இருந்திருக்கும்}}.) Did it happen? (\textbf{No.})
\end{tcolorbox}
\section[(six lines)]{Six lines about a trip}

\label{lesson:TA-W293-six-lines}



\begin{quote}
The last connected message you wrote was about a familiar day. This one tells a story: six lines about a trip, to your home town or anywhere else.

Pick one reader: a friend, by name.
\end{quote}

\subsection*{What to know first}
A told story is not a list of events. It has a \textbf{scene}: \textbf{-\ta{கொண்டிருந்தேன்}} or \textbf{-\ta{த} \ta{போது}}. It has \textbf{events} in the plain past, kept in order with \textbf{\ta{முதலில்}}, \textbf{\ta{அதன்} \ta{பிறகு}}, the \textbf{-\ta{உ}} / \textbf{-\ta{இ}} form and \textbf{\ta{கடைசியில்}}. It has at least one \textbf{reason}, with \textbf{\ta{அதனால்}}, \textbf{-\ta{ததால்}} or \textbf{\ta{ஏனென்றால்}}, and a \textbf{feeling}, with \textbf{\ta{எனக்கு} … \ta{வந்தது}} or \textbf{\ta{எனக்கு} … -\ta{ஆக} \ta{இருந்தது}}.

A good one ends by looking at what might have been: \textbf{… -\ta{இருந்தால்}, … -\ta{இருப்பேன்}}.

\subsection*{Writing — connected composition}
On paper, in Tamil script, write \textbf{six lines} to your reader about a trip, real or invented, one line for each of these:

\begin{enumerate}
\raggedright
  \item the scene: when it was, and what you were doing as it began, with \textbf{-\ta{கொண்டிருந்தேன்}} or \textbf{-\ta{த} \ta{போது}};
  \item the first thing you did, with \textbf{\ta{முதலில்}}: a ticket, a seat, or the way you asked;
  \item what happened next, with \textbf{\ta{அதன்} \ta{பிறகு}} or an \textbf{-\ta{உ}} / \textbf{-\ta{இ}} form, and one thing that went wrong: a missed bus, a missing bag, or rain and a wet shirt;
  \item why, or what followed, with \textbf{\ta{அதனால்}}, \textbf{-\ta{ததால்}} or \textbf{\ta{ஏனென்றால்}};
  \item how it ended and how you felt, with \textbf{\ta{கடைசியில்}} and \textbf{\ta{எனக்கு} … \ta{வந்தது}} or \textbf{-\ta{ஆக} \ta{இருந்தது}};
  \item one line about what might have been, with \textbf{-\ta{இருந்தால்} … -\ta{இருப்பேன்}}.
\end{enumerate}

There is no model on this page, and no word bank. Use only words you already write. The whole exercise is choosing the right past for each line, and joining the lines.

\begin{tcolorbox}[breakable,colback=teal!4,colframe=teal!35!black,title={Before you move on}]
Now read your lines aloud. Mark every verb \textbf{s} for scene and \textbf{e} for event. Line one should be mostly \textbf{s}; lines two and three mostly \textbf{e}.

Then a spelling check. If it rained in your story, did you write \textbf{\ta{ஈரம்}} or \textbf{\ta{ஈரமாக}}? It opens with the letter \textbf{\ta{ஈ}}: a bar, two legs, two dots, not the curled \textbf{\ta{இ}}. Inside a word, after a consonant, the same long \emph{ī} is the sign \ta{ீ}.

Last, read line six alone. Does the if-clause end in \textbf{-\ta{இருந்தால்}}, and the main verb in \textbf{-\ta{இருப்பேன்}}?
\end{tcolorbox}
\section[(review)]{Review — it depends, and what would have been}

\label{lesson:TA-R293-depends-and-would}



\begin{quote}
Say \textquotedblleft{}It depends on the weather\textquotedblright{} (\emph{atu vāṉilaiyaip poṟuttatu}) and \textquotedblleft{}if (you) had come\textquotedblright{} (\emph{vantiruntāl}).

\begin{itemize}
\raggedright
  \item \emph{Write it:} \textbf{\ta{ஈ}} and \textbf{\ta{ஈரம்}}, once each
\end{itemize}
\end{quote}

\subsection*{Your turn}
Say these aloud:

\begin{itemize}
\raggedright
  \item Will you come tomorrow? It depends on the rain. If it rains from morning until evening, I will stay at home.
  \item If you don't come soon, the bus will leave.
  \item If I had had time, I would have come.
  \item If you had come, we would have gone to the sea together.
  \item read your six lines aloud to your reader
\end{itemize}

\begin{tcolorbox}[breakable,colback=teal!4,colframe=teal!35!black,title={Before you move on}]
\textquotedblleft{}It depends on …\textquotedblright{}? (\textbf{… -\ta{ஐப்} \ta{பொறுத்தது}}.) \textquotedblleft{}If (you) had come\textquotedblright{}? (\textbf{\ta{வந்திருந்தால்}}.) \textquotedblleft{}I would have come\textquotedblright{}? (\textbf{\ta{வந்திருப்பேன்}}.) \textquotedblleft{}Depending on the price\textquotedblright{}? (\textbf{\ta{விலையைப்} \ta{பொறுத்து}}.)
\end{tcolorbox}
